%%%PAGE 23 rot=0 legibility=good summary=电磁感应单杆模型：过程/状态/动能三种分析方法(冲量、牛二、动能定理与能量)、电热在串并联电阻中的分配、微积分处理变力冲量、动量定理推导I_A=±BLq=±B²L²x/R=m(v2-v1)
%%%BLOCK id=023-01 ch=12 sec=单杆模型·分析方法 kind=method alt=none cont=none y=0.075-0.20
\textbf{单杆}　过程/状态/动能
\[
\left\{\begin{aligned}
F_A\ \text{恒力}\quad&W_A=\pm BILx\\
F_A\ \text{变力}\quad&\pm W_A
\end{aligned}\right.
\]
冲量：
\begin{align*}
I_A&=BIL\cdot\Delta t\quad\text{过程写动\unclear{能}动量}\\
I_A&=\pm BLq
\end{align*}
% CHECK: “过程写动能/动量”处“能”与“量”字迹重叠，疑为把“动能”改成“动量”（动量定理），照原样
状态写牛二　最大　最小　稳定$\Leftrightarrow a=0$
%%%END
%%%BLOCK id=023-02 ch=12 sec=电磁感应·能量 kind=knowledge alt=none cont=none y=0.20-0.31
\begin{itemize}
  \item 功和功关系　动能定理
  \item 功和能关系　动能定理
% CHECK: 原稿第二行右侧同为“动能定理”（“功”字有蓝色荧光笔涂抹），疑应为“功能关系”，照原样
  \item 能和能关系　能量守恒
  \begin{itemize}
    \item 机械能　机械能定理
  \end{itemize}
\end{itemize}
%%%END
%%%BLOCK id=023-03 ch=12 sec=电磁感应·电热分配 kind=method alt=10 cont=none y=0.31-0.62
\textbf{电热/功/功率　分配}
%FIG p023_a 0.075 0.365 0.215 0.442
\notefig[0.3]{figures/p023_a.png}
%FIG p023_b 0.23 0.362 0.35 0.435
\notefig[0.3]{figures/p023_b.png}
\begin{center}
\begin{tabular}{p{0.44\linewidth}p{0.44\linewidth}}
并联　$Q$与$R$成反比 & 串联　$Q$与$R$成正比\\[2pt]
$\dfrac{U^2}{R}t$ & $I^2Rt$\\[6pt]
$Q_1=\dfrac{R_2}{R_1+R_2}Q$ & $Q_1=\dfrac{R_1}{R_1+R_2}Q$\\[8pt]
$Q_2=\dfrac{R_1}{R_1+R_2}Q$ & $Q_2=Q\dfrac{R_2}{R_1+R_2}$\\[8pt]
$Q$和其它电阻成正比 & $Q$和自身电阻成正比
\end{tabular}
\end{center}
$P$、$W$ 一样
%%%END
%%%BLOCK id=023-04 ch=12 sec=电磁感应·变力冲量(微积分) kind=derivation alt=90 cont=none y=0.62-0.72
\textbf{微积分}
\[
\left\{\begin{aligned}
&f=kv\\
&qvB\\
&\frac{B^2L^2v}{R}
\end{aligned}\right.
\qquad a\text{变化的运动}
\]
% CHECK: “a变化的运动”首字形似 Q/a，按“a（加速度）变化的运动”转录
\[
I=\int f\,\dd t=k\,v\,\dd t=kx
\]
% CHECK: 原稿第二个等号后为“k vdt”，无积分号，照原样
%%%END
%%%BLOCK id=023-05 ch=12 sec=电磁感应·动量 kind=derivation alt=none cont=none y=0.72-0.99
\[
I_A=\pm BLq=\pm\frac{B^2L^2x}{R}=m(v_2-v_1)
\]
（$BLq$ 下方标注：）流过导体电荷量\\
（箭头 $\uparrow$ 指向 $x$，标注：）单杆导棒位移/双杆位移差
\begin{align*}
-BIL\,\Delta t&=m\,\Delta v\\
-BL\,\Delta q&=m\,\Delta v\\
-\frac{B^2L^2v}{R_{\text{总}}}\,\Delta t&=m\,\Delta v\\
-\frac{B^2L^2x}{R_{\text{总}}}&=m\,\Delta v
\end{align*}
求和：
\begin{align*}
-BLq&=m(v_2-v_1)\\
-\frac{B^2L^2x}{R_{\text{总}}}&=m(v_2-v_1)
\end{align*}
%%%END
%%%PAGEEND 23
